\section{משוואות דיפרנציאליות חלקיות אי־הומוגניות — סיכום }

\begin{summarybox}{מטרות הלמידה}
בסיום שיעור זה תדעו:
\begin{itemize}
    \item לזהות מד"ח אי־הומוגנית ולהבחין בין אי־הומוגניות במשוואה לבין אי־הומוגניות בתנאי השפה
    \item להשתמש בשיטת פריסה לפונקציות עצמיות לפתרון בעיות כלליות
    \item למצוא פתרון פרטי/סטציונרי ולהפחית בעיה אי־הומוגנית לבעיה הומוגנית
    \item להבחין בין התנהגות דיפוזיה (דעיכה) לבין גלים (תנודות מתמשכות)
\end{itemize}
\end{summarybox}

\subsection{הבעיה הכללית — הגדרות ומוטיבציה}

\begin{definitionbox}{מד"ח אי־הומוגנית}
משוואה דיפרנציאלית חלקית נקראת \textbf{אי־הומוגנית} אם היא מכילה איבר שאינו תלוי בפונקציה הנעלמת $u$ או בנגזרותיה. איבר זה נקרא \textbf{כוח מאלץ} או \textbf{מקור}.
\end{definitionbox}

נבחן שתי דוגמאות מרכזיות:

\textbf{משוואת גלים עם כוח מאלץ} (למשל, מיתר תחת כבידה או כוח חיצוני):
\[
\boxed{\frac{\partial^2 u}{\partial t^2}-c^2\frac{\partial^2 u}{\partial x^2}=F(x,t)}
\]

\textbf{משוואת דיפוזיה עם מקור חום/חומר} (למשל, מוט עם חימום פנימי):
\[
\boxed{\frac{\partial u}{\partial t}=a^2\frac{\partial^2 u}{\partial x^2}+F(x,t)}
\]

הכוח $F$ עשוי להיות תלוי בזמן ובמקום, ותנאי השפה וההתחלה יכולים להיות הומוגניים או אי־הומוגניים.

\begin{notebox}{עקרון הסופרפוזיציה — הבסיס לכל השיטות}
עבור משוואות \textbf{ליניאריות}, הפתרון הכללי הוא תמיד:
\[
u_{\text{כללי}} = u_{\text{הומוגני}} + u_{\text{פרטי}}
\]
זוהי הבסיס לכל השיטות שנלמד בשיעור זה.
\end{notebox}

\subsection{שיטה א': פירוק לטורי פונקציות עצמיות (השיטה הכללית)}

\subsubsection{הרעיון המרכזי}
נניח תנאי שפה הומוגניים $u(0,t)=u(L,t)=0$. הפונקציות העצמיות של האופרטור $-\frac{d^2}{dx^2}$ עם תנאים אלו הן:
\[
\phi_n(x) = \sin\!\left(\frac{n\pi x}{L}\right), \qquad n=1,2,3,\dots
\]
עם ערכים עצמיים $\lambda_n = \left(\frac{n\pi}{L}\right)^2$.

\subsubsection{שלב 1: פריסת הכוח המאלץ}
נפרוס את הכוח המאלץ בבסיס הפונקציות העצמיות:
\[
F(x,t)=\sum_{n=1}^\infty F_n(t)\sin\!\left(\frac{n\pi x}{L}\right)
\]

\subsubsection{שלב 2: חישוב המקדמים $F_n(t)$}
המקדמים $F_n(t)$ מחושבים באמצעות אורתוגונליות הפונקציות העצמיות:
\[
\boxed{F_n(t)=\frac{2}{L}\int_0^L F(x,t)\sin\!\left(\frac{n\pi x}{L}\right)\,dx}
\]

\textbf{הוכחה מפורטת:} נכפיל את הפריסה של $F$ ב־$\sin\left(\frac{m\pi x}{L}\right)$ ונאינטגרל על $[0,L]$:
\begin{align*}
\int_0^L F(x,t)\sin\!\left(\frac{m\pi x}{L}\right)dx &= \int_0^L \sum_{n=1}^\infty F_n(t) \sin\!\left(\frac{n\pi x}{L}\right)\sin\!\left(\frac{m\pi x}{L}\right)dx \\
&= \sum_{n=1}^\infty F_n(t) \int_0^L \sin\!\left(\frac{n\pi x}{L}\right)\sin\!\left(\frac{m\pi x}{L}\right)dx
\end{align*}
נשתמש ביחס האורתוגונליות:
\[
\int_0^L \sin\!\left(\frac{n\pi x}{L}\right)\sin\!\left(\frac{m\pi x}{L}\right)dx = \frac{L}{2}\delta_{nm} = \begin{cases} \frac{L}{2} & n=m \\ 0 & n \neq m \end{cases}
\]
לכן:
\[
\int_0^L F(x,t)\sin\!\left(\frac{m\pi x}{L}\right)dx = F_m(t) \cdot \frac{L}{2} \implies F_m(t) = \frac{2}{L}\int_0^L F(x,t)\sin\!\left(\frac{m\pi x}{L}\right)dx
\]

\subsubsection{שלב 3: הנחת צורת הפתרון}
נניח פתרון מהצורה:
\[
u(x,t)=\sum_{n=1}^\infty h_n(t)\sin\!\left(\frac{n\pi x}{L}\right)
\]
המקדמים $h_n(t)$ הם פונקציות לא ידועות של הזמן שיש למצוא.

\subsubsection{שלב 4: הצבה במשוואת הגלים — גזירה מלאה}
נציב את הפריסות במשוואה $\frac{\partial^2 u}{\partial t^2}-c^2\frac{\partial^2 u}{\partial x^2}=F(x,t)$:

\textbf{נגזרת שנייה בזמן:}
\[
\frac{\partial^2 u}{\partial t^2} = \frac{\partial^2}{\partial t^2}\sum_{n=1}^\infty h_n(t)\sin\!\left(\frac{n\pi x}{L}\right) = \sum_{n=1}^\infty \ddot{h}_n(t)\sin\!\left(\frac{n\pi x}{L}\right)
\]

\textbf{נגזרת שנייה במרחב:}
\[
\frac{\partial^2 u}{\partial x^2} = \sum_{n=1}^\infty h_n(t) \cdot \frac{\partial^2}{\partial x^2}\sin\!\left(\frac{n\pi x}{L}\right) = \sum_{n=1}^\infty h_n(t) \cdot \left(-\frac{n^2\pi^2}{L^2}\right)\sin\!\left(\frac{n\pi x}{L}\right)
\]

\textbf{הצבה במשוואה המלאה:}
\[
\sum_{n=1}^\infty \ddot{h}_n(t)\sin\!\left(\frac{n\pi x}{L}\right) + c^2\sum_{n=1}^\infty h_n(t)\frac{n^2\pi^2}{L^2}\sin\!\left(\frac{n\pi x}{L}\right) = \sum_{n=1}^\infty F_n(t)\sin\!\left(\frac{n\pi x}{L}\right)
\]

\[
\sum_{n=1}^\infty \left[\ddot{h}_n(t) + \frac{c^2 n^2\pi^2}{L^2}h_n(t)\right]\sin\!\left(\frac{n\pi x}{L}\right) = \sum_{n=1}^\infty F_n(t)\sin\!\left(\frac{n\pi x}{L}\right)
\]

\textbf{השוואת מקדמים:} מאחר שהפונקציות $\sin\left(\frac{n\pi x}{L}\right)$ בלתי תלויות ליניארית, נקבל לכל $n$:
\[
\boxed{\ddot{h}_n(t) + \omega_n^2 h_n(t) = F_n(t), \qquad \omega_n = \frac{n\pi c}{L}}
\]

זוהי משוואת \textbf{אוסילטור הרמוני מאולץ} — משוואה דיפרנציאלית רגילה מסדר שני!

\subsubsection{שלב 5: פתרון המד"ר עבור כל מוד}
הפתרון הכללי למשוואה $\ddot{h}_n + \omega_n^2 h_n = F_n(t)$ הוא סכום של פתרון הומוגני ופתרון פרטי:
\[
h_n(t)=h_n^{(h)}(t)+h_n^{(p)}(t)
\]

\textbf{הפתרון ההומוגני} (אוסילטור הרמוני חופשי):
\[
h_n^{(h)}(t)=A_n\cos(\omega_n t)+B_n\sin(\omega_n t)
\]

\textbf{הפתרון הפרטי} באמצעות וריאציית פרמטרים (\LR{Variation of Parameters}):

נזכר בנוסחה הכללית: עבור משוואה $y'' + \omega^2 y = g(t)$, הפתרון הפרטי הוא:
\[
y_p(t) = \frac{1}{\omega}\int_0^t g(\tau)\sin\bigl(\omega(t-\tau)\bigr)\,d\tau
\]

לכן:
\[
\boxed{h_n^{(p)}(t) = \frac{1}{\omega_n}\int_0^t F_n(\tau)\sin\bigl(\omega_n(t-\tau)\bigr)\,d\tau}
\]

\begin{notebox}{השפעת תנאי ההתחלה — נקודה קריטית!}
המקדמים $A_n, B_n$ נקבעים מתנאי ההתחלה $u(x,0)=u_0(x)$ ו־$\partial_t u(x,0)=v_0(x)$. 

\textbf{חשוב:} יש להציב את הפתרון \textbf{המלא} (כולל הפתרון הפרטי!) בתנאי ההתחלה:
\begin{align*}
u(x,0) &= \sum_{n=1}^\infty \left[A_n + h_n^{(p)}(0)\right]\sin\!\left(\frac{n\pi x}{L}\right) = u_0(x) \\
\partial_t u(x,0) &= \sum_{n=1}^\infty \left[B_n\omega_n + \dot{h}_n^{(p)}(0)\right]\sin\!\left(\frac{n\pi x}{L}\right) = v_0(x)
\end{align*}
\end{notebox}

\subsection{שיטה ב': הפרדת פתרון פרטי/סטציונרי (השיטה היעילה)}

כאשר $F$ פשוט (למשל תלוי רק ב־$x$ ולא בזמן), קיימת שיטה קצרה יותר — נוח לנחש פתרון פרטי $u_p(x)$ שאינו תלוי בזמן:
\[
u(x,t)=u_h(x,t)+u_p(x)
\]
אם $u_p$ מקיים את המשוואה המלאה, אז $u_h$ חייב לקיים את המשוואה ההומוגנית, אך עם תנאי התחלה מתוקנים (הזזה):
\[
u_h(x,0)=u_0(x)-u_p(x), \qquad
\partial_t u_h(x,0)=v_0(x)
\]

\subsubsection{גזירת המשוואה עבור $u_h$ — הוכחה מלאה}
נציב $u = u_p + u_h$ במשוואת הגלים:
\begin{align*}
\frac{\partial^2 (u_p + u_h)}{\partial t^2} - c^2\frac{\partial^2 (u_p + u_h)}{\partial x^2} &= F(x) \\
\underbrace{\frac{\partial^2 u_h}{\partial t^2}}_{u_p \text{ לא תלוי ב-}t} - c^2\left(\frac{d^2 u_p}{dx^2} + \frac{\partial^2 u_h}{\partial x^2}\right) &= F(x) \\
\frac{\partial^2 u_h}{\partial t^2} - c^2\frac{\partial^2 u_h}{\partial x^2} &= F(x) + c^2 \frac{d^2 u_p}{dx^2}
\end{align*}

אם בחרנו $u_p$ כך ש־$-c^2 u_p'' = F(x)$, אז:
\[
\frac{\partial^2 u_h}{\partial t^2} - c^2\frac{\partial^2 u_h}{\partial x^2} = 0
\]
כלומר $u_h$ מקיים משוואה \textbf{הומוגנית}!

\subsection{טיפול בתנאי שפה אי־הומוגניים}

\begin{problembox}{הבעיה}
אם תנאי השפה אינם אפס, למשל:
\[
u(0,t)=U_1, \qquad u(L,t)=U_2
\]
לא ניתן להשתמש ישירות בטורי סינוסים (כי הם מתאפסים בקצוות).
\end{problembox}

\subsubsection{הפתרון — שיטת ההזזה}

\textbf{שלב 1: מציאת פונקציית הזזה}

נחפש פונקציה פשוטה $u_s(x)$ המקיימת את תנאי השפה. הבחירה הפשוטה ביותר היא פונקציה ליניארית:
\[
u_s(x) = ax + b
\]
נדרוש:
\begin{itemize}
    \item $u_s(0) = b = U_1 \Rightarrow b = U_1$
    \item $u_s(L) = aL + U_1 = U_2 \Rightarrow a = \frac{U_2 - U_1}{L}$
\end{itemize}
לכן:
\[
\boxed{u_s(x) = U_1 + \frac{U_2 - U_1}{L}x}
\]

\textbf{שלב 2: הגדרת משתנה חדש}
\[
v(x,t) = u(x,t) - u_s(x)
\]

\textbf{שלב 3: בדיקת התנאים עבור $v$}
\begin{itemize}
    \item תנאי שפה: $v(0,t) = u(0,t) - u_s(0) = U_1 - U_1 = 0$ \checkmark
    \item תנאי שפה: $v(L,t) = u(L,t) - u_s(L) = U_2 - U_2 = 0$ \checkmark
\end{itemize}
הפונקציה $v$ מקיימת תנאי שפה \textbf{הומוגניים}!

\textbf{שלב 4: המשוואה עבור $v$}

המשוואה עבור $v$ עשויה להפוך לאי-הומוגנית אם $u_s$ לא פותר את המשוואה המקורית. עכשיו ניתן לפתור עבור $v$ בשיטות הרגילות.
מאוד חשוב לזכור:
ב-T=0
\[v(x,t=0)=\overbrace{g(x)}^{u(x,t=0)}-v_{st}(x)\]
קיבלנו תנאי התחלה מתוקן.
הפתרון יהיה:
\[u(x,t)=v(x,t)+u_{st}(x,t)=\underbrace{\sum a_n \sin(\frac{\pi nx}{L}) e^{-\frac{a^2 \pi^2 n^2}{L^2}}}_v+\underbrace{U_1 + \frac{U_2 - U_1}{L}x}_u \]


\subsection{משוואת הדיפוזיה — ניתוח מקיף}

\subsubsection{התנהגות אסימפטוטית}
במשוואת הדיפוזיה, המודים הגבוהים דועכים מהר (כי $e^{-a^2 n^2 \pi^2 t / L^2}$ קטן יותר עבור $n$ גדול). לכן, לאחר זמן רב ($t \to \infty$), המערכת שואפת ל\textbf{מצב סטציונרי} (אם קיים).

\subsubsection{מקרה א': קיים מצב סטציונרי}
אם קיים פתרון סטציונרי $u_s(x)$ כך ש־$\partial_t u_s=0$, אזי מהמשוואה $\partial_t u = a^2 \partial_{xx} u + F$:
\[
0 = a^2\frac{d^2 u_s}{dx^2}+F(x) \implies \boxed{\frac{d^2 u_s}{dx^2} = -\frac{F(x)}{a^2}}
\]

\textbf{פתרון באינטגרציה כפולה:}
\begin{align*}
\frac{du_s}{dx} &= -\frac{1}{a^2}\int F(x)\,dx + C_1 \\
u_s(x) &= -\frac{1}{a^2}\iint F(x)\,dx\,dx + C_1 x + C_2
\end{align*}
הקבועים $C_1, C_2$ נקבעים מתנאי השפה.

\textbf{דוגמה ללא מקור:} ללא כוח חיצוני ($F=0$) ובתנאי שפה $u(0)=U_1$, $u(L)=U_2$:
\[
\frac{d^2 u_s}{dx^2} = 0 \implies u_s = C_1 x + C_2
\]
מתנאי השפה: $C_2 = U_1$ ו־$C_1 = \frac{U_2-U_1}{L}$, לכן:
\[
u_s(x)=U_1+\frac{U_2-U_1}{L}x
\]
זהו פרופיל ליניארי!

\subsubsection{הפתרון הכללי לדיפוזיה}
הפתרון הכללי הוא סכום של המצב הסטציונרי וטרנזיינט שדועך בזמן:
\[
\boxed{u(x,t)=u_s(x)+v(x,t)}
\]
כאשר הדעיכה היא אקספוננציאלית:
\[
v(x,t)=\sum_{n=1}^\infty A_n \sin\!\left(\frac{n\pi x}{L}\right) e^{-a^2(n\pi/L)^2 t}
\]
המקדמים $A_n$ נקבעים מתנאי ההתחלה המתוקן: $v(x,0) = u_0(x) - u_s(x)$.



\subsubsection{מקרה ב': אין מצב סטציונרי (הצטברות אנרגיה)}
אם תנאי השפה הם של שטף (\LR{Neumann}) כך שסך האנרגיה הנכנסת שונה מהיוצאת — אין פתרון סטציונרי. נחפש פתרון פרטי מהצורה $u_p(x,t) = At + \phi(x)$. דוגמה מפורטת מופיעה בסעיף הבא.

\subsection{משוואת גלים — הבדלים מהותיים מדיפוזיה}

כאשר $F(x)$ אינו תלוי בזמן (למשל מיתר תחת כבידה), ניתן למצוא \emph{צורת ייחוס} (שווי משקל) $u_p(x)$ כך ש:
\[
-c^2u_p''(x)=F(x) \implies \boxed{u_p''(x) = -\frac{F(x)}{c^2}}
\]
הפתרון הכללי מתאר תנודות סביב צורת שווי המשקל החדשה:
\[
u(x,t)=u_p(x)+u_h(x,t)
\]

\begin{resultbox}{ההבדל המהותי בין גלים לדיפוזיה}
\begin{center}
\renewcommand{\arraystretch}{1.5}
\begin{tabular}{|c|c|c|}
\hline
& \textbf{דיפוזיה} & \textbf{גלים} \\
\hline
התנהגות בזמן ארוך & דעיכה למצב סטציונרי & תנודות נצחיות \\
\hline
משמעות $u_p$ & מצב סופי של המערכת & צורת שיווי משקל סביבה מתנדנדים \\
\hline
קיום פתרון סטציונרי & מתקיים לרוב & תמיד מתנדנד \\
\hline
\end{tabular}
\end{center}

במשוואת גלים (ללא חיכוך) אין דעיכה ל"מצב סטציונרי". הפתרון הפרטי $u_p(x)$ מגדיר את נקודת שיווי המשקל שסביבה המיתר מתנודד \textbf{לנצח}, בעוד שבדיפוזיה המערכת קורסת לתוך הפתרון הסטציונרי כשהזמן שואף לאינסוף.
\end{resultbox}

\begin{examplebox}{דוגמה: מיתר תחת כבידה}
\textbf{נתון:}
\[
\frac{\partial^2 u}{\partial t^2} - c^2\frac{\partial^2 u}{\partial x^2} = -g, \qquad u(0,t)=u(L,t)=0
\]

\textbf{שלב 1: מציאת צורת שיווי משקל}
\[
-c^2\frac{d^2 u_p}{dx^2} = -g \implies \frac{d^2 u_p}{dx^2} = \frac{g}{c^2}
\]

\textbf{שלב 2: אינטגרציה ראשונה}
\[
\frac{du_p}{dx} = \frac{g}{c^2}x + C_1
\]

\textbf{שלב 3: אינטגרציה שנייה}
\[
u_p(x) = \frac{g}{2c^2}x^2 + C_1 x + C_2
\]

\textbf{שלב 4: קביעת קבועים מתנאי שפה}
\begin{itemize}
    \item $u_p(0) = 0 \Rightarrow C_2 = 0$
    \item $u_p(L) = 0 \Rightarrow \frac{g}{2c^2}L^2 + C_1 L = 0 \Rightarrow C_1 = -\frac{gL}{2c^2}$
\end{itemize}

\textbf{צורת שיווי המשקל:}
\[
u_p(x) = \frac{g}{2c^2}x^2 - \frac{gL}{2c^2}x = \frac{g}{2c^2}x(x-L)
\]
\[
\boxed{u_p(x) = \frac{g}{2c^2}x(x-L)}
\]
זוהי \textbf{פרבולה הפוכה} — המיתר "שוקע" באמצע!

\textbf{שלב 5: הפתרון המלא}
\[
u(x,t) = u_p(x) + \sum_{n=1}^\infty \left[A_n\cos(\omega_n t) + B_n\sin(\omega_n t)\right]\sin\!\left(\frac{n\pi x}{L}\right)
\]
כאשר $\omega_n = \frac{n\pi c}{L}$.

\textbf{תנאי התחלה מתוקנים:}
\begin{itemize}
    \item מיקום: $v(x,0) = u_0(x) - u_p(x)$ (משתנה!)
    \item מהירות: $\partial_t v(x,0) = v_0(x)$ (לא משתנה!)
\end{itemize}

המיתר מתנדנד \textbf{לנצח} סביב צורת השיווי משקל $u_p(x)$.
\end{examplebox}



\newpage
\section{משוואת פואסון ונושאים נוספים}

בחלק זה נרחיב את הדיון לכלול את משוואת פואסון בקואורדינטות פולריות.

\subsection{דוגמאות נוספות למשוואת הדיפוזיה}

נציג שתי דוגמאות מפורטות הממחישות את השיטות שלמדנו.

\begin{examplebox}{דוגמה מפורטת: מוט עם מקור חום קבוע $Q$}
עבור $F(x) = Q$ (כאשר $Q$ קבוע) ותנאי שפה $u(0)=U_1$, $u(L)=U_2$:

\textbf{שלב 1:} הפתרון הסטציונרי מקיים $\frac{d^2 u_s}{dx^2} = -\frac{Q}{a^2}$.

\textbf{שלב 2:} אינטגרציה כפולה נותנת פרבולה:
\begin{equation}
u_s(x) = -\frac{Q}{2a^2}x^2 + C_1 x + C_2
\end{equation}

\textbf{שלב 3:} מתנאי שפה: $C_2 = U_1$ ו־$C_1 = \frac{U_2-U_1}{L} + \frac{QL}{2a^2}$.

\textbf{התוצאה:}
\[
\boxed{u_s(x) = -\frac{Q}{2a^2}x^2 + \left(\frac{U_2-U_1}{L}+\frac{QL}{2a^2}\right)x + U_1}
\]

הפתרון המלא הוא $u = u_s + v$, כאשר $v$ דועך לאפס בזמן אינסוף.
\end{examplebox}

\begin{figure}[h]
\centering
\begin{tikzpicture}
\begin{axis}[
width=8cm, height=5cm,
xlabel={$x$}, ylabel={$u$},
xmin=0, xmax=1, ymin=0, ymax=1.5,
xtick={0,1}, xticklabels={$0$, $L$},
ytick={0.5, 1}, yticklabels={$U_1$, $U_2$},
axis lines=middle,
title={התכנסות למצב סטציונרי (פרבולה) בדיפוזיה}
]
% Steady state
\addplot[blue, thick, domain=0:1] {-0.8*x^2 + 1.3*x + 0.5};
\node[blue] at (axis cs: 0.5, 1.1) {$u_s(x)$};

% Transient states
\addplot[red, dashed, domain=0:1] {0.5 + 0.5*x}; % Initial could be anything
\addplot[gray, domain=0:1] {-0.4*x^2 + 0.9*x + 0.5}; % Intermediate

\end{axis}
\end{tikzpicture}
\caption{התפתחות הפתרון מתנאי התחלה שרירותי (קו מקווקו) דרך מצבי ביניים (אפור) אל המצב הסטציונרי (כחול).}
\end{figure}

\begin{examplebox}{דוגמה: מוט מבודד בצד אחד ושטף קבוע בשני (ללא מצב סטציונרי)}
\textbf{נתון:} $u_x(0,t) = 0$ (בידוד), $u_x(L,t) = B^2$ (שטף נכנס), ללא מקור פנימי ($F=0$).

מכיוון שחום נכנס ואין לו לאן לצאת, המערכת מתחממת לינארית.

\textbf{ניחוש:} $u_p(x,t) = At + \frac{B^2}{2L}x^2$.

\textbf{בדיקה:}
\begin{itemize}
\item הצבה במשוואה: $A = a^2 \cdot \frac{B^2}{L} \Rightarrow A = \frac{a^2 B^2}{L}$
\item תנאי שפה ב-$x=0$: $\partial_x u_p = \frac{B^2}{L} \cdot 0 = 0$ \checkmark
\item תנאי שפה ב-$x=L$: $\partial_x u_p = \frac{B^2}{L} \cdot L = B^2$ \checkmark
\end{itemize}

\textbf{התוצאה:}
\[
\boxed{u_p(x,t) = \frac{a^2 B^2}{L}t + \frac{B^2}{2L}x^2}
\]
כך קיבלנו פתרון פרטי המקיים את המשוואה ותנאי השפה הלא-הומוגניים.
\end{examplebox}

\subsection{משוואת פואסון במעגל}

משוואת פואסון היא משוואת לפלס אי-הומוגנית: $\nabla^2 u = \rho$.
נפתור דוגמה בקואורדינטות פולריות בתוך עיגול $r < R$ עם צפיפות מטען אחידה $\rho = 1$.

\begin{examplebox}{דוגמה מפורטת: צפיפות מטען אחידה בדיסק}
\textbf{הבעיה:} פתור $\nabla^2 u = 1$ בתוך דיסק ברדיוס $R$.

\textbf{שלב 1 — לפלסיאן בקואורדינטות פולריות:}

בסימטריה מעגלית ($u$ לא תלוי ב־$\theta$):
\begin{equation}
\nabla^2 u = \frac{1}{r}\frac{d}{dr}\left(r\frac{du}{dr}\right) = 1
\end{equation}

\textbf{שלב 2 — ניחוש פתרון פרטי $u_p(r)$:}

נכפיל ב־$r$:
\[
\frac{d}{dr}\left(r\frac{du_p}{dr}\right) = r
\]

\textbf{שלב 3 — אינטגרציה ראשונה:}
\[
r\frac{du_p}{dr} = \frac{r^2}{2} + C_1 \implies \frac{du_p}{dr} = \frac{r}{2} + \frac{C_1}{r}
\]

\textbf{שלב 4 — אינטגרציה שנייה:}
\begin{equation}
u_p(r) = \frac{r^2}{4} + C_1 \ln r + C_2
\end{equation}

\textbf{שלב 5 — תנאים פיזיקליים:}

כדי למנוע התבדרות בראשית $r = 0$, נזכר ש־$\ln r \to -\infty$ כש־$r \to 0$.
לכן נדרוש:
\[
\boxed{C_1 = 0}
\]

את $C_2$ ניתן לבחור כרצוננו (למשל כדי לאפס את הפתרון על השפה $r=R$), או להשאיר ולתקן בפתרון ההומוגני.

\textbf{הפתרון הפרטי:}
\[
\boxed{u_p(r) = \frac{r^2}{4} + C_2}
\]

\textbf{הפתרון המלא:}
\[
u(r,\theta) = u_p(r) + u_h(r,\theta)
\]
כאשר $u_h$ פותר את משוואת לפלס ההומוגנית $\nabla^2 u_h = 0$ ומקיים את תנאי השפה המתוקנים על $r = R$.
\end{examplebox}

\subsection{טבלת סיכום  — בחירת שיטת הפתרון}

\begin{theorembox}{סיכום שיטות — מדריך מהיר}
\begin{center}
\renewcommand{\arraystretch}{1.6}
\begin{tabular}{|p{4.5cm}|p{7cm}|}
\hline
\textbf{סוג הבעיה} & \textbf{שיטה מומלצת} \\
\hline
\textbf{כוח תלוי זמן ומקום} $F(x,t)$ & שיטת הפריסה לפונקציות עצמיות (שטורם-ליוביל) \\
\hline
\textbf{כוח תלוי מקום בלבד (דיפוזיה)} & מציאת פתרון סטציונרי $u_s(x)$ ופתרון הבעיה ההומוגנית עבור ההפרש $v = u - u_s$ \\
\hline
\textbf{כוח תלוי מקום בלבד (גלים)} & מציאת צורת שיווי משקל $u_p(x)$ ותנודות $v$ סביבה \\
\hline
\textbf{ללא שיווי משקל (דיפוזיה)} & הפרדת משתנים מוכללת: $u_p(x,t) = At + \phi(x)$ \\
\hline
\textbf{תנאי שפה אי־הומוגניים} & הזזה: $v = u - u_s$ כאשר $u_s$ (לרוב ליניארי) מקיים את תנאי השפה \\
\hline
\textbf{סימטריה מעגלית (פואסון)} & ניחוש $u_p(r)$, דרישת רגולריות בראשית ($C_1 = 0$) \\
\hline
\end{tabular}
\end{center}
\end{theorembox}

\begin{summarybox}{עקרונות מנחים }
\begin{enumerate}
    \item \textbf{עקרון הסופרפוזיציה:} כל בעיה אי־הומוגנית ליניארית ניתנת לפירוק:
    \[
    u = u_{\text{פרטי}} + u_{\text{הומוגני}}
    \]
    
    \item \textbf{דיפוזיה:} דעיכה למצב סטציונרי (אם קיים). קצב הדעיכה של מוד $n$:
    \[
    e^{-a^2 n^2 \pi^2 t/L^2}
    \]
    
    \item \textbf{גלים:} תנודות נצחיות סביב שיווי משקל. אין דעיכה ללא מנגנון חיכוך.
    
    \item \textbf{תנאי שפה אי־הומוגניים:} הזזה לתנאים הומוגניים:
    \[
    v = u - u_s, \qquad u_s(0) = U_1, \quad u_s(L) = U_2
    \]
    
    \item \textbf{תנאי התחלה:} משתנים בהתאם לפתרון הפרטי — תמיד לזכור לתקן!
    \[
    v(x,0) = u_0(x) - u_p(x)
    \]
    
    \item \textbf{רגולריות:} בבעיות עם סימטריה (דיסק, כדור), לדרוש התנהגות סופית בראשית.
\end{enumerate}
\end{summarybox}